\documentclass[10pt,letter]{report}

\usepackage{import}
\import{../../../../sistema/}{rutas}
\usepackage[T1]{fontenc}
%\usepackage[letterpaper, headsep=60pt, headheight=2cm]{geometry}
\usepackage{verbatim}
\usepackage[utf8x]{inputenc}
\usepackage[table]{xcolor}
\usepackage{float}
\usepackage{import}
\usepackage{textcomp}
\usepackage{ifthen}
\usepackage[spanish,mexico-com]{babel}
\usepackage{pdflscape}
\usepackage{mathrsfs}
\usepackage{amsmath}
\usepackage{amssymb}
\usepackage{bbm}
\usepackage{tikz}
\usepackage{fp}
\usepackage[autolanguage]{numprint}
\usepackage{array}
\usetikzlibrary{shapes}
\usepackage{subfigure}
\usepackage{ucs}
\usepackage[utf8x]{inputenc}
\usepackage{fontenc}
\usepackage{graphicx}
\usepackage{anysize}
\usepackage{relsize}
\usepackage{booktabs}

\usepackage{fancyhdr}
\usepackage[all]{xy}
\setlength{\headheight}{13.1pt}
\makeatletter\renewcommand\theenumi{\@alph\c@enumi}\makeatother
\renewcommand\labelenumi{\theenumi)}
\usepackage{amsthm}
\usepackage{enumerate}

\usepackage[]{mdframed}

\usepackage{marvosym}
\usepackage{tikzsymbols}
\usepackage{imakeidx}%for indexes
%\usepackage{tocbibind}
\usepackage{background}
\usepackage{titlesec}
\usepackage{multirow}
\usepackage{etoolbox}
\usepackage{fmtcount}
\usepackage{datetime}
\usepackage[bookmarks=true]{hyperref}%must be at the end of preamble
\usepackage[toc,section=subsection, acronym, shortcuts]{glossaries}
\usepackage[open,openlevel=0]{bookmark}
\usepackage{lipsum}
\usepackage{multicol}
\usepackage{wrapfig}
\usepackage[titletoc]{appendix}
\usepackage{sectsty}

\usepackage{titlesec}
\usepackage{pdfpages}
%---------------Formato Numeros \numprint-------------
\npdecimalsign{.}
\nprounddigits{2}
\npthousandsep{,}

%-----------------Espacio en blanco--------------------
\newcommand{\espacio}[1]{\vspace{#1}\begin{center}\textit{[EL RESTO DE ESTA PÁGINA HA SIDO DEJADO EN BLANCO\\
DE FORMA INTENCIONAL]}\end{center}\newpage}
\newcommand{\inserta}{\colorbox{principal}{\textcolor{orange}{Insertar}}}
\usepackage{textcomp}
\decimalpoint
%------------Dise�o de p�gina --------------------------------
\usepackage[centering, letterpaper,margin=2cm,top=1cm, headsep=24pt, headheight=2cm,includehead, includefoot]{geometry}

%------------Colores----------------------

\definecolor{principal}{RGB}{0, 53, 73}
\definecolor{secundario}{RGB}{43, 82, 96}
\definecolor{terciario}{RGB}{43,82,96}

%--------------Hipervinculos en color, ligas-azul, archivos-magenta, url-azul----------------
\hypersetup{
    colorlinks=true,
    linkcolor=secundario,
    filecolor=magenta,      
    urlcolor=blue}

%-------------Encabezado---------------------
\pagestyle{fancy}
\fancyhf{}

\chead{\begin{tabular}{m{14.2cm}m{2.5cm}}
\rowcolor{principal}
\includegraphics[width=5cm]{\rutaImagenes/logovaluaminew.pdf} &\includegraphics[width=2.5cm]{\rutaImagenes/isotipo_valuami_oscuro}\\
\end{tabular}
}

%--------------Pie de p�gina------------------
\cfoot{\textbf{\textcolor{principal}{\footnotesize{VALUAMI\tiny\textregistered}}}\\ \scriptsize{\textit{Montes Urales 755 Piso 5, Col. Lomas de Chapultepec, Miguel Hidalgo,\\ CDMX, C.P. 11000. Tel. (55)41616060 ext.248 / Cel. (55) 54 17 18 46.\\Calle 1 n\'umero 233, Local 25
Colonia Arboledas,\\ Villahermosa Tabasco, C.P. 86079, Tel. (993) 186 39 95}}}
\rfoot{\footnotesize{\thepage\hspace{3pt}de \pageref{lastpage}}}
\lfoot{
\begin{tabular}{b{2cm}}
\includegraphics[width=1.5cm]{\rutaImagenes/qr}
\end{tabular}}

%------------T�tulos de seccion y subseccion--------------



\renewcommand \thechapter {\Roman{chapter}}
\renewcommand \thesection {\Alph{section}}
\renewcommand \thesubsection {\thesection.\arabic{subsection}}
\renewcommand \thesubsubsection {\thesection.\arabic{subsection}.\arabic{subsubsection}}

\titleformat{\section}[hang]{\color{gray}\huge\bfseries}{\thesection.}{1em}{} 

\chapterfont{\color{principal}}
\sectionfont{\color{principal}}
\subsectionfont{\color{secundario}}
\subsubsectionfont{\color{terciario}}

\titleformat{\chapter}[display]{\normalfont\Large\filcenter\sffamily}
{\titlerule[1pt]%
\vspace{1pt}%
\titlerule
\vspace{1pc}%
\LARGE\MakeUppercase{\chaptertitlename} \thechapter}
{1pc}
{\titlerule
\vspace{1pc}%
\Huge}
\titlespacing{\chapter}{0pt}{*-8}{*1.5}

%------------------Profundidad del �ndice---------------------------

\setcounter{tocdepth}{3}
\setcounter{secnumdepth}{3}


%------------------Marca de agua---------------

\backgroundsetup{angle=0, contents={\includegraphics[width=4cm]{\rutaImagenes/isotipo_valuami.pdf}},opacity=.2, scale=5}


%----------------------------Generales----------------------------
\newcommand{\tipoAvaluo}{insertar}
\newcommand{\bienesValuados}{insertar}%nombre del bien que se va a valuar
\newcommand{\lugarInforme}{insertar}

%--------------------Fechas---------------------------
\newcommand{\diainforme}{insertar} %dia del informe
\newcommand{\mesinforme}{insertar} %mes del informe
\newcommand{\annoinforme}{insertar} %año del informe

\newcommand{\diavalores}{insertar} %dia de valores
\newcommand{\mesvalores}{insertar} %mes de valores
\newcommand{\annovalores}{insertar} %año de valores

\newcommand{\diainspeccion}{insertar} %dia de inspeccion
\newcommand{\mesinspeccion}{insertar} %mes de inspeccion
\newcommand{\annoinspeccion}{insertar} %año de inspeccion

\newcommand{\fechaInforme}{\diainforme{} de \monthname[\mesinforme] de \annoinforme}
\newcommand{\fechaValores}{\diavalores{} de \monthname[\mesvalores] de \annovalores}
\newcommand{\fechaValoresCorto}{\diavalores/\mesvalores/\annovalores}
\newcommand{\fechaInspeccion}{\diainspeccion{} de \monthname[\mesinspeccion] de \annoinspeccion}

%----------------------------Perito Valuador---------------------
\newcommand{\peritoValuador}{insertar}
\import{\rutaValuatex/documentos_modelo/peritos/peritos/}{\peritoValuador}

%------------------------Perito Auxiliar-------------------
\newcommand{\peritoAuxiliar}{insertar}
\ifthenelse{\equal{\peritoAuxiliar}{n/a}}{ }{\import{\rutaValuatex/documentos_modelo/peritos/auxiliares/}{\peritoAuxiliar}}

%--------------Datos del solicitante-------------
\newcommand{\empresaSolicitante}{insertar}
\newcommand{\empresaCorto}{insertar}
\newcommand{\rfcEmpresa}{insertar}
\newcommand{\personaSolicitante}{insertar}
\newcommand{\caracterSolicitante}{insertar}

%--------------Datos del propietario-------------
\newcommand{\nombrePropietario}{insertar}

%-------------Vigencia----------------------------
\newcommand{\vigenciaInforme}{insertar}
\newcommand{\notaVigencia}{si}


%---------------Ubicacion del bien sujeto de valuaci\'on
\newcommand{\descripcionBien}{insertar}%en que consiste el bien que se va a valuar.
\newcommand{\ubicacionBien}{insertar}
\newcommand{\rfcBien}{insertar}

%------------Uso de la valuación---------------
\newcommand{\usoAvaluo}{insertar}
%--------------Desarrollo del avalúo en la especie-------------------
\newcommand{\EFde}{insertar}
\newcommand{\EFhasta}{insertar}
\newcommand{\EFdeHasta}{insertar}

%----------------Parámetros Peers--------------------------
%__________Multiplo________________
\newcommand{\peersa}{insertar}
\newcommand{\peersb}{insertar}
\newcommand{\peersc}{insertar}
\newcommand{\peersd}{insertar}
\newcommand{\peerse}{insertar}

%____________x veces________
\newcommand{\peersaTo}{insertar}
\newcommand{\peersbTo}{insertar}
\newcommand{\peerscTo}{insertar}
\newcommand{\peersdTo}{insertar}
\newcommand{\peerseTo}{insertar}

%___________Valor Múltiplo____________
\newcommand{\peersaMult}{insertar}
\newcommand{\peersbMult}{insertar}
\newcommand{\peerscMult}{insertar}
\newcommand{\peersdMult}{insertar}
\newcommand{\peerseMult}{insertar}

%___________Estadístico____________

\newcommand{\peersaEst}{insertar}
\newcommand{\peersbEst}{insertar}
\newcommand{\peerscEst}{insertar}
\newcommand{\peersdEst}{insertar}
\newcommand{\peerseEst}{insertar}

%============ WACC ===============
%---------------------- RF ------------------------------
\newcommand{\rfBase}{insertar}
\newcommand{\rfValor}{insertar}

%---------------------- Beta ----------------------------
\newcommand{\valorBeta}{insertar}
%-----------------------Premio de Mercado (ERP)----------------------------
\newcommand{\mercadoAccionario}{insertar}%mexicano o americano
\newcommand{\erpValor}{insertar}

%----------------------Riesgo Pais (CRP)-----------------
\newcommand{\crpValor}{insertar}

%-----------------------Costo de Capital (ke)--------------------
\newcommand{\keValor}{insertar}

%----------------------Costo de deuda (Kd)-----------------------
\newcommand{\kdValor}{insertar}

%----------------------- Valor Wacc----------------------------
\newcommand{\waccValor}{insertar}

%=============== DCF ======================
\newcommand{\periodoProyeccion}{insertar}
\newcommand{\proyCagr}{insertar}
\newcommand{\proyEbitMargin}{insertar}

\newcommand{\tasaFiscal}{insertar}
\newcommand{\reinvestmentRate}{insertar}


%============== RFR (Relief from royalties)==========
\newcommand{\tasaRegalias}{insertar}
\newcommand{\estadisticoTasaRegalias}{insertar}


%============== CIFRAS ======================
\newcommand{\valorDCF}{insertar}
\newcommand{\valorDCFLetra}{insertar} 

\newcommand{\valorPEERS}{insertar}
\newcommand{\valorPEERSLetra}{insertar} 

\newcommand{\valorFirma}{insertar}
\newcommand{\valorFirmaLetra}{insertar} 

\newcommand{\valorCapital}{insertar}
\newcommand{\valorCapitalLetra}{insertar} 

\newcommand{\valorRFR}{insertar}
\newcommand{\valorRFRLetra}{insertar} 

\newcommand{\valorResidual}{insertar}
\newcommand{\valorResidualLetra}{insertar} 

\newcommand{\valorActivoIntangible}{insertar}
\newcommand{\valorActivoIntangibleLetra}{insertar} 

\newcommand{\valorCapitalIntangible}{insertar}
\newcommand{\valorCApitalIntangibleLetra}{insertar} 

\newcommand{\moneda}{insertar}
\newcommand{\monedaCode}{insertar}









\makenoidxglossaries

\newglossaryentry{beta}
{
        name=Beta,
        description={El coeficiente Beta ($\beta$)} es una medida de volatilidad que estima el riesgo sistem\'atico de un activo.
}

\newglossaryentry{cashNeq}
{
        name=Cash \& Eq,
        description={Efectivo y equivalentes}
}

\newglossaryentry{efectivetaxrate}
{
        name=Effective Tax Rate,
        description={Tasa fiscal efectiva}
}

\newglossaryentry{enterprisevalue}
{
        name=Enterprise value,
        description={Valor de la empresa}
}

\newglossaryentry{equityvalue}
{
        name=Equity value,
        description={Valor del capital accionario}
}

\newglossaryentry{exitmultiple}
{
        name=Exit multiple,
        description={M\'ultiplo de salida}
}

\newglossaryentry{expdebt}
{
        name=Exp. Debt,
        description={Deuda expl\'icita}
}

\newglossaryentry{fairvalue}
{
        name=Fair Value,
        description={Valor razonable o valor justo de mercado}
}

\newglossaryentry{firmvalue}
{
        name=Firm Value,
        description={Valor de la firma o empresa}
}

\newglossaryentry{growthrate}
{
        name=Growth Rate,
        description={Tasa de crecimiento de ingresos netos}
}

\newglossaryentry{incomestatement}
{
        name=Income Statement,
        description={Estado de Resultados}
}

\newglossaryentry{investmentcapital}
{
        name=Investment Capital,
        description={Capital Invertido (IC)}
}

\newglossaryentry{leveredbeta}
{
        name=Levered beta,
        description={Beta apalancada}
}

\newglossaryentry{marketaproach}
{
        name=Market approach,
        description={Enfoque de mercado}
}

\newglossaryentry{netdebt}
{
        name=Net Debt,
        description={Deuda neta}
}

\newglossaryentry{nopat}
{
        name=Nopat,
        description={Flujo de Operaci\'on Neto}
}

\newglossaryentry{projectvaluation}
{
        name=Project Valuation,
        description={Valor razonable del proyecto de inversi\'on}
}

\newglossaryentry{riskfreerate}
{
        name=Risk Free Rate,
        description={Tasa Libre de Riesgo}
}

\newglossaryentry{sizeprime}
{
        name=Size Prime,
        description={Prima por tama\~no}
}

\newglossaryentry{terminalvalue}
{
        name=Terminal Value,
        description={Valor terminal}
}

\newglossaryentry{totaldebt}
{
        name=Total Debt,
        description={Deuda total}
}

\newglossaryentry{valuedrivers}
{
        name=Value drivers,
        description={Elevadores de Valor}
}
\newacronym{cagr}{CAGR}{Tasa compuesta de crecimiento anual}
\newacronym{capm}{CAPM}{Modelo de Valoraci\'on de Activos de Capital}
\newacronym{chic}{CHIC}{Sociedades de Inversi\'on Cerradas}
\newacronym{crp}{CRP}{Riesgo Pa\'is}

\newacronym{dcf}{DCF}{Flujo de efectivo descontado (Discounted Cash FLow)}
\newacronym{erp}{ERP}{Premio de Mercado}
\newacronym{etr}{ETR}{Tasa Fiscal Efectiva}
\newacronym{fcff}{FCFF}{Flujo de efectivo libre a la Firma (Free Cash Flow to Firm)}
\newacronym{fcfe}{FCFE}{Flujo de efectivo libre al Capital (Free Cash Flow to Equity)}
\newacronym{g}{G}{Tasa de Crecimiento de ingresos netos (Grow Rate)}
\newacronym{ke}{Ke}{Costo de Capital Accionario}
\newacronym{kd}{Kd}{Costo de la deuda}
\newacronym{kpi}{KPI}{Indicador clave de desempe\~no o indicadores de gesti\'on \textit{(Key Performance Indicator)}}
\newacronym{mpeem}{MPEEM}{Multi-Period Excess Earnings Method}
\newacronym{nav}{NAV}{Valor neto de activos}
\newacronym{nca}{NCA}{Activo No Corriente}
\newacronym{nwc}{NWC}{Capital de Trabajo Neto}
\newacronym{peers}{PEERS}{M\'ultiplos de Cotizaci\'on}
\newacronym{rfr}{RFR}{M\'etodo de Flujo de Ahorro en Regalias}
\newacronym{wacc}{WACC}{Costo Promedio Ponderado de Capital}
\newacronym{wara}{WARA}{Costo Promedio Ponderado de Activos}

\usepackage{pdfpages}

\begin{document}
\def\sectionautorefname{Inciso}
\def\subsectionautorefname{Inciso}

\fechaInforme


\chapter{Uno}\label{1}
\thispagestyle{fancy}
\section{uno}\label{1.1}
\subsection{1}\label{1.1.1}
\subsubsection{1}

Hola\autoref{1}, \autoref{1.1}, \autoref{1.1.1}
\spanishdate{2021}{12}{9} \... \ldots  1000.00, $10000.00$
\newpage
`a', "<a">, ``a'' \gls{beta}, \glspl{leveredbeta} \acs{cagr} \acs{capm}, \acs{dcf}

en  \numberstringnum{4} \monthname[1] a los \numberstringnum{\diainforme} d\'ias del mes de \monthname[\mesinforme] del a\~no \numberstringnum{\annoinforme}
\\
\diainforme{} d\'ias del mes de \monthname[\mesinforme]{} del a\~no \annoinforme (\numberstringnum{\annoinforme})
\begin{center}
 \printnoidxglossary[type=\acronymtype,title={Acr\'onimos}]
\printnoidxglossary
\end{center}

\begin{table}[H]
\centering
\begin{tabular}{|c|m{.4\textwidth}|}
\hline
Obra& Dibujo	\\
\hline
Autor & Se atribuye a Ra\'ul Gonzalez Esquivel\\
\hline
T\'itulo & ``Figura S\'edente, Cultura Totonaca, Veracruz, Cl\'asico Tard\'io, Septiembre 29 de 1993'' \\
\hline
Tem\'atica & Representaci\'on de arte precolombino, im\'agenes de deidades, figuras antropomorfas, cabezas colosales, culto religioso ind\'igena y personajes mitol\'ogicos.\\
\hline
T\'ecnicas &Grafito sobre papel \\
\hline
Medidas & 41.5 cm $\times$ 59 cm \\
\hline
 Formato & Peque\~no \\
 \hline
 Firma & Se atribuye a Ra\'ul Gonzalez Esquivel \\
 \hline
 Fecha & \\
 \hline
 Provenance & Colecci\'on particular\\
 \hline
 Certificado de autenticidad& Firmado por Pablo C. Goebel el 10 de abril de 1996.  \\
 \hline 
  Materiales empleados & Lapiz sobre papel\\
 \hline
 Caracter\'isticas especiales & Forma parte de una colecci\'on privada de 39 dibujos. 
Contiene dibujo con el texto ``Bufete Campos'' que corresponde a la obra por encargo. \\
\hline 
Ubicaci\'on & Calle Constituci\'on 124, colonia Centro, Veracruz, Veracruz, C.P. 91700.\\
\hline
\end{tabular}
\end{table}

\begin{table}[H]
\centering
\begin{tabular}{|c|m{.4\textwidth}|}
\hline
Obra& Dibujo	\\
\hline
Autor & Se atribuye a Ra\'ul Gonzalez Esquivel\\
\hline
T\'itulo & ``Dios Viejo (Remojadas, Estado de Veracruz). Cultura Totonaca. Septiembre 12 de 1993'' \\
\hline
Tem\'atica & Representaci\'on de arte precolombino, im\'agenes de deidades, figuras antropomorfas, cabezas colosales, culto religioso ind\'igena y personajes mitol\'ogicos.\\
\hline
T\'ecnicas &Grafito sobre papel \\
\hline
Medidas & 41.5 cm $\times$ 59 cm \\
\hline
 Formato & Peque\~no \\
 \hline
 Firma & Se atribuye a Ra\'ul Gonzalez Esquivel \\
 \hline
  Fecha & \\
 \hline
 Provenance & Colecci\'on particular\\
 \hline
 Certificado de autenticidad& Firmado por Pablo C. Goebel el 10 de abril de 1996.  \\
 \hline 
  Materiales empleados & Lapiz sobre papel\\
 \hline
 Caracter\'isticas especiales & Forma parte de una colecci\'on privada de 39 dibujos. 
Contiene dibujo con el texto ``Bufete Campos'' que corresponde a la obra por encargo. \\
\hline 
Ubicaci\'on & Calle Constituci\'on 124, colonia Centro, Veracruz, Veracruz, C.P. 91700.\\
\hline

\end{tabular}
\end{table}

\begin{table}[H]
\centering
\begin{tabular}{|c|m{.4\textwidth}|}
\hline
Obra& Dibujo	\\
\hline
Autor & Se atribuye a Ra\'ul Gonzalez Esquivel\\
\hline
T\'itulo & ``Representaci\'on de Quetzalcoatl, Taj\'in, Ve. Cultura Totonaca. Septiembre 20 de 1993'.'' \\
\hline
Tem\'atica & Representaci\'on de arte precolombino, im\'agenes de deidades, figuras antropomorfas, cabezas colosales, culto religioso ind\'igena y personajes mitol\'ogicos.\\
\hline
T\'ecnicas &Grafito sobre papel \\
\hline
Medidas & 41.5 cm $\times$ 59 cm \\
\hline
 Formato & Peque\~no \\
 \hline
 Firma & Se atribuye a Ra\'ul Gonzalez Esquivel  \\
 \hline
 Fecha & \\
 \hline
 Provenance & Colecci\'on particular\\
 \hline
 Certificado de autenticidad& Firmado por Pablo C. Goebel el 10 de abril de 1996.  \\
 \hline 
  Materiales empleados & Lapiz sobre papel\\
 \hline
 Caracter\'isticas especiales & Forma parte de una colecci\'on privada de 39 dibujos. 
Contiene dibujo con el texto ``Bufete Campos'' que corresponde a la obra por encargo. \\
\hline 
Ubicaci\'on & Calle Constituci\'on 124, colonia Centro, Veracruz, Veracruz, C.P. 91700.\\
\hline

\end{tabular}
\end{table}

\begin{table}[H]
\centering
\begin{tabular}{|c|m{.4\textwidth}|}
\hline
Obra& Dibujo	\\
\hline
Autor & Se atribuye a Ra\'ul Gonzalez Esquivel\\
\hline
T\'itulo & ``Nichos del edificio no. 16 del Taj\'in. El Taj\'in Veracruz. Septiembre 5 de 1993.''\\
\hline
Tem\'atica & Representaci\'on de arte precolombino, im\'agenes de deidades, figuras antropomorfas, cabezas colosales, culto religioso ind\'igena y personajes mitol\'ogicos.\\
\hline
T\'ecnicas &Grafito sobre papel \\
\hline
Medidas & 41.5 cm $\times$ 59 cm \\
\hline
 Formato & Peque\~no \\
 \hline
 Firma & Se atribuye a Ra\'ul Gonzalez Esquivel  \\
 \hline
  Fecha & \\
 \hline
 Provenance & Colecci\'on particular\\
 \hline
 Certificado de autenticidad& Firmado por Pablo C. Goebel el 10 de abril de 1996.  \\
 \hline 
  Materiales empleados & Lapiz sobre papel\\
 \hline
 Caracter\'isticas especiales & Forma parte de una colecci\'on privada de 39 dibujos. 
Contiene dibujo con el texto ``Bufete Campos'' que corresponde a la obra por encargo. \\
\hline 
Ubicaci\'on & Calle Constituci\'on 124, colonia Centro, Veracruz, Veracruz, C.P. 91700.\\
\hline

\end{tabular}
\end{table}

\begin{table}[H]
\centering
\begin{tabular}{|c|m{.4\textwidth}|}
\hline
Obra& Dibujo	\\
\hline
Autor & Se atribuye a Ra\'ul Gonzalez Esquivel\\
\hline
T\'itulo & ``Cabeza de Anciano. Tres Zapotes, La Venta, Veracruz. Septiembre 11 de 1993.'' \\
\hline
Tem\'atica & Representaci\'on de arte precolombino, im\'agenes de deidades, figuras antropomorfas, cabezas colosales, culto religioso ind\'igena y personajes mitol\'ogicos.\\
\hline
T\'ecnicas &Grafito sobre papel \\
\hline
Medidas & 41.5 cm $\times$ 59 cm \\
\hline
 Formato & Peque\~no \\
 \hline
 Firma & Se atribuye a Ra\'ul Gonzalez Esquivel \\
 \hline
  Fecha & \\
 \hline
 Provenance & Colecci\'on particular\\
 \hline
 Certificado de autenticidad& Firmado por Pablo C. Goebel el 10 de abril de 1996.  \\
 \hline 
  Materiales empleados & Lapiz sobre papel\\
 \hline
 Caracter\'isticas especiales & Forma parte de una colecci\'on privada de 39 dibujos. 
Contiene dibujo con el texto ``Bufete Campos'' que corresponde a la obra por encargo. \\
\hline 
Ubicaci\'on & Calle Constituci\'on 124, colonia Centro, Veracruz, Veracruz, C.P. 91700.\\
\hline

\end{tabular}
\end{table}

\begin{table}[H]
\centering
\begin{tabular}{|c|m{.4\textwidth}|}
\hline
Obra& Dibujo	\\
\hline
Autor & Se atribuye a Ra\'ul Gonzalez Esquivel\\
\hline
T\'itulo & `Èscultura Masculina S\'edente. El Zapotal, Mpio. de Ignacio de la Llave. S. VI-IX. Agosto 23 de 1993.'' \\
\hline
Tem\'atica & Representaci\'on de arte precolombino, im\'agenes de deidades, figuras antropomorfas, cabezas colosales, culto religioso ind\'igena y personajes mitol\'ogicos.\\
\hline
T\'ecnicas &Grafito sobre papel \\
\hline
Medidas & 41.5 cm $\times$ 59 cm \\
\hline
 Formato & Peque\~no \\
 \hline
 Firma &  Se atribuye a Ra\'ul Gonzalez Esquivel\\
 \hline
  Fecha & \\
 \hline
 Provenance & Colecci\'on particular\\
 \hline
 Certificado de autenticidad& Firmado por Pablo C. Goebel el 10 de abril de 1996.  \\
 \hline 
  Materiales empleados & Lapiz sobre papel\\
 \hline
 Caracter\'isticas especiales & Forma parte de una colecci\'on privada de 39 dibujos. 
Contiene dibujo con el texto ``Bufete Campos'' que corresponde a la obra por encargo. \\
\hline 
Ubicaci\'on & Calle Constituci\'on 124, colonia Centro, Veracruz, Veracruz, C.P. 91700.\\
\hline

\end{tabular}
\end{table}

\begin{table}[H]
\centering
\begin{tabular}{|c|m{.4\textwidth}|}
\hline
Obra& Dibujo	\\
\hline
Autor & Se atribuye a Ra\'ul Gonzalez Esquivel\\
\hline
T\'itulo & ``Mensajero del Sol. El Zapotal, Mpio. de Ignacio de la Llave. Cl\'asico Tard\'io. S. VI-IX D.C. Septiembre 03 de 1993.''\\
\hline
Tem\'atica & Representaci\'on de arte precolombino, im\'agenes de deidades, figuras antropomorfas, cabezas colosales, culto religioso ind\'igena y personajes mitol\'ogicos.\\
\hline
T\'ecnicas &Grafito sobre papel \\
\hline
Medidas & 41.5 cm $\times$ 59 cm \\
\hline
 Formato & Peque\~no \\
 \hline
 Firma & Se atribuye a Ra\'ul Gonzalez Esquivel \\
 \hline
  Fecha & \\
 \hline
 Provenance & Colecci\'on particular\\
 \hline
 Certificado de autenticidad& Firmado por Pablo C. Goebel el 10 de abril de 1996.  \\
 \hline 
  Materiales empleados & Lapiz sobre papel\\
 \hline
 Caracter\'isticas especiales & Forma parte de una colecci\'on privada de 39 dibujos. 
Contiene dibujo con el texto ``Bufete Campos'' que corresponde a la obra por encargo. \\
\hline 
Ubicaci\'on & Calle Constituci\'on 124, colonia Centro, Veracruz, Veracruz, C.P. 91700.\\
\hline

\end{tabular}
\end{table}

\begin{table}[H]
\centering
\begin{tabular}{|c|m{.4\textwidth}|}
\hline
Obra& Dibujo	\\
\hline
Autor & Se atribuye a Ra\'ul Gonzalez Esquivel\\
\hline
T\'itulo & ``Sin t\'itulo, del Arq. Gonz\'alea Esquivel, aunque si muestra padre con hijo en sus brazos, o bien personaje con persona menor en brazos. Agosto 14 de 1993'' \\
\hline
Tem\'atica & Representaci\'on de arte precolombino, im\'agenes de deidades, figuras antropomorfas, cabezas colosales, culto religioso ind\'igena y personajes mitol\'ogicos.\\
\hline
T\'ecnicas &Grafito sobre papel \\
\hline
Medidas & 41.5 cm $\times$ 59 cm \\
\hline
 Formato & Peque\~no \\
 \hline
 Firma & Se atribuye a Ra\'ul Gonzalez Esquivel\\ 
 \hline
  Fecha &  Agosto 14 de 1993\\
 \hline
 Provenance & Colecci\'on particular\\
 \hline
 Certificado de autenticidad& Firmado por Pablo C. Goebel el 10 de abril de 1996.  \\
 \hline 
  Materiales empleados & Lapiz sobre papel\\
 \hline
 Caracter\'isticas especiales & Forma parte de una colecci\'on privada de 39 dibujos. 
Contiene dibujo con el texto ``Bufete Campos'' que corresponde a la obra por encargo. \\
\hline 
Ubicaci\'on & Calle Constituci\'on 124, colonia Centro, Veracruz, Veracruz, C.P. 91700.\\
\hline

\end{tabular}
\end{table}

\begin{table}[H]
\centering
\begin{tabular}{|c|m{.4\textwidth}|}
\hline
Obra& Dibujo	\\
\hline
Autor & Se atribuye a Ra\'ul Gonzalez Esquivel\\
\hline
T\'itulo & ``Dos cabezas, hombre y mujer, paga, Agosto 22 de 1993''\\
\hline
Tem\'atica & Representaci\'on de arte precolombino, im\'agenes de deidades, figuras antropomorfas, cabezas colosales, culto religioso ind\'igena y personajes mitol\'ogicos.\\
\hline
T\'ecnicas &Grafito sobre papel \\
\hline
Medidas & 41.5 cm $\times$ 59 cm \\
\hline
 Formato & Peque\~no \\
 \hline
 Firma & Se atribuye a Ra\'ul Gonzalez Esquivel\\ 
 \hline
  Fecha & Agosto 22 de 1993\\
 \hline
 Provenance & Colecci\'on particular\\
 \hline
 Certificado de autenticidad& Firmado por Pablo C. Goebel el 10 de abril de 1996.  \\
 \hline 
  Materiales empleados & Lapiz sobre papel\\
 \hline
 Caracter\'isticas especiales & Forma parte de una colecci\'on privada de 39 dibujos. 
Contiene dibujo con el texto ``Bufete Campos'' que corresponde a la obra por encargo. \\
\hline 
Ubicaci\'on & Calle Constituci\'on 124, colonia Centro, Veracruz, Veracruz, C.P. 91700.\\
\hline

\end{tabular}
\end{table}

\begin{table}[H]
\centering
\begin{tabular}{|c|m{.4\textwidth}|}
\hline
Obra& Dibujo	\\
\hline
Autor & Se atribuye a Ra\'ul Gonzalez Esquivel\\
\hline
T\'itulo & ``Tlalsolteotl, protectora de los partos y proporcionadora de los amores il\'icitos y las pasiones. Agosto de 1993.''\\
\hline
Tem\'atica & Representaci\'on de arte precolombino, im\'agenes de deidades, figuras antropomorfas, cabezas colosales, culto religioso ind\'igena y personajes mitol\'ogicos.\\
\hline
T\'ecnicas &Grafito sobre papel \\
\hline
Medidas & 41.5 cm $\times$ 59 cm \\
\hline
 Formato & Peque\~no \\
 \hline
 Firma & Se atribuye a Ra\'ul Gonzalez Esquivel\\ 
 \hline
  Fecha & Agosto de 1993.\\
 \hline
 Provenance & Colecci\'on particular\\
 \hline
 Certificado de autenticidad& Firmado por Pablo C. Goebel el 10 de abril de 1996.  \\
 \hline 
  Materiales empleados & Lapiz sobre papel\\
 \hline
 Caracter\'isticas especiales & Forma parte de una colecci\'on privada de 39 dibujos. 
Contiene dibujo con el texto ``Bufete Campos'' que corresponde a la obra por encargo. \\
\hline 
Ubicaci\'on & Calle Constituci\'on 124, colonia Centro, Veracruz, Veracruz, C.P. 91700.\\
\hline

\end{tabular}
\end{table}

\begin{table}[H]
\centering
\begin{tabular}{|c|m{.4\textwidth}|}
\hline
Obra& Dibujo	\\
\hline
Autor & Se atribuye a Ra\'ul Gonzalez Esquivel\\
\hline
T\'itulo & ``Dios Taj\'in, Dios de la lluvia y del trueno, edificio no. 5. Taj\'in Papantla, Veracruz. Septiembre de 1993.''\\
\hline
Tem\'atica & Representaci\'on de arte precolombino, im\'agenes de deidades, figuras antropomorfas, cabezas colosales, culto religioso ind\'igena y personajes mitol\'ogicos.\\
\hline
T\'ecnicas &Grafito sobre papel \\
\hline
Medidas & 41.5 cm $\times$ 59 cm \\
\hline
 Formato & Peque\~no \\
 \hline
 Firma & Se atribuye a Ra\'ul Gonzalez Esquivel\\ 
 \hline
  Fecha & Septiembre de 1993.\\
 \hline
 Provenance & Colecci\'on particular\\
 \hline
 Certificado de autenticidad& Firmado por Pablo C. Goebel el 10 de abril de 1996.  \\
 \hline 
  Materiales empleados & Lapiz sobre papel\\
 \hline
 Caracter\'isticas especiales & Forma parte de una colecci\'on privada de 39 dibujos. 
Contiene dibujo con el texto ``Bufete Campos'' que corresponde a la obra por encargo. \\
\hline 
Ubicaci\'on & Calle Constituci\'on 124, colonia Centro, Veracruz, Veracruz, C.P. 91700.\\
\hline

\end{tabular}
\end{table}

\begin{table}[H]
\centering
\begin{tabular}{|c|m{.4\textwidth}|}
\hline
Obra& Dibujo	\\
\hline
Autor & Se atribuye a Ra\'ul Gonzalez Esquivel\\
\hline
T\'itulo &``Carita SOnriente de los cerros. Cultura Totonaca (Veracruz). Mayo 15 de 1993.'' \\
\hline
Tem\'atica & Representaci\'on de arte precolombino, im\'agenes de deidades, figuras antropomorfas, cabezas colosales, culto religioso ind\'igena y personajes mitol\'ogicos.\\
\hline
T\'ecnicas &Grafito sobre papel \\
\hline
Medidas & 41.5 cm $\times$ 59 cm \\
\hline
 Formato & Peque\~no \\
 \hline
 Firma & Se atribuye a Ra\'ul Gonzalez Esquivel\\ 
 \hline
  Fecha & Mayo 15 de 1993.\\
 \hline
 Provenance & Colecci\'on particular\\
 \hline
 Certificado de autenticidad& Firmado por Pablo C. Goebel el 10 de abril de 1996.  \\
 \hline 
  Materiales empleados & Lapiz sobre papel\\
 \hline
 Caracter\'isticas especiales & Forma parte de una colecci\'on privada de 39 dibujos. 
Contiene dibujo con el texto ``Bufete Campos'' que corresponde a la obra por encargo. \\
\hline 
Ubicaci\'on & Calle Constituci\'on 124, colonia Centro, Veracruz, Veracruz, C.P. 91700.\\
\hline

\end{tabular}
\end{table}

\begin{table}[H]
\centering
\begin{tabular}{|c|m{.4\textwidth}|}
\hline
Obra& Dibujo	\\
\hline
Autor & Se atribuye a Ra\'ul Gonzalez Esquivel\\
\hline
T\'itulo & ``Carita Sonriente. Cultura Totonaca. Julio 19 de 1993.'' \\
\hline
Tem\'atica & Representaci\'on de arte precolombino, im\'agenes de deidades, figuras antropomorfas, cabezas colosales, culto religioso ind\'igena y personajes mitol\'ogicos.\\
\hline
T\'ecnicas &Grafito sobre papel \\
\hline
Medidas & 41.5 cm $\times$ 59 cm \\
\hline
 Formato & Peque\~no \\
 \hline
 Firma & Se atribuye a Ra\'ul Gonzalez Esquivel\\ 
 \hline
  Fecha & Julio 19 de 1993.\\
 \hline
 Provenance & Colecci\'on particular\\
 \hline
 Certificado de autenticidad& Firmado por Pablo C. Goebel el 10 de abril de 1996.  \\
 \hline 
  Materiales empleados & Lapiz sobre papel\\
 \hline
 Caracter\'isticas especiales & Forma parte de una colecci\'on privada de 39 dibujos. 
Contiene dibujo con el texto ``Bufete Campos'' que corresponde a la obra por encargo. \\
\hline 
Ubicaci\'on & Calle Constituci\'on 124, colonia Centro, Veracruz, Veracruz, C.P. 91700.\\
\hline

\end{tabular}
\end{table}

\begin{table}[H]
\centering
\begin{tabular}{|c|m{.4\textwidth}|}
\hline
Obra& Dibujo	\\
\hline
Autor & Se atribuye a Ra\'ul Gonzalez Esquivel\\
\hline
T\'itulo & `Hacha Votiva con Mono, Centro de Veracruz Post-Cl\'asico Tard\'io S. VI-IX D.C.''\\
\hline
Tem\'atica & Representaci\'on de arte precolombino, im\'agenes de deidades, figuras antropomorfas, cabezas colosales, culto religioso ind\'igena y personajes mitol\'ogicos.\\
\hline
T\'ecnicas &Grafito sobre papel \\
\hline
Medidas & 41.5 cm $\times$ 59 cm \\
\hline
 Formato & Peque\~no \\
 \hline
 Firma & Se atribuye a Ra\'ul Gonzalez Esquivel\\ 
 \hline
  Fecha & 6 de diciembre de 1993.\\
 \hline
 Provenance & Colecci\'on particular\\
 \hline
 Certificado de autenticidad& Firmado por Pablo C. Goebel el 10 de abril de 1996.  \\
 \hline 
  Materiales empleados & Lapiz sobre papel\\
 \hline
 Caracter\'isticas especiales & Forma parte de una colecci\'on privada de 39 dibujos. 
Contiene dibujo con el texto ``Bufete Campos'' que corresponde a la obra por encargo. \\
\hline 
Ubicaci\'on & Calle Constituci\'on 124, colonia Centro, Veracruz, Veracruz, C.P. 91700.\\
\hline

\end{tabular}
\end{table}

\begin{table}[H]
\centering
\begin{tabular}{|c|m{.4\textwidth}|}
\hline
Obra& Dibujo	\\
\hline
Autor & Se atribuye a Ra\'ul Gonzalez Esquivel\\
\hline
T\'itulo & ``Carita Sonriente de los Cerros. Cultura Totonaca (Veracruz). Mayo 8 de 1993.''\\
\hline
Tem\'atica & Representaci\'on de arte precolombino, im\'agenes de deidades, figuras antropomorfas, cabezas colosales, culto religioso ind\'igena y personajes mitol\'ogicos.\\
\hline
T\'ecnicas &Grafito sobre papel \\
\hline
Medidas & 41.5 cm $\times$ 59 cm \\
\hline
 Formato & Peque\~no \\
 \hline
 Firma & Se atribuye a Ra\'ul Gonzalez Esquivel\\ 
 \hline
  Fecha & Mayo 8 de 1993.\\
 \hline
 Provenance & Colecci\'on particular\\
 \hline
 Certificado de autenticidad& Firmado por Pablo C. Goebel el 10 de abril de 1996.  \\
 \hline 
  Materiales empleados & Lapiz sobre papel\\
 \hline
 Caracter\'isticas especiales & Forma parte de una colecci\'on privada de 39 dibujos. 
Contiene dibujo con el texto ``Bufete Campos'' que corresponde a la obra por encargo. \\
\hline 
Ubicaci\'on & Calle Constituci\'on 124, colonia Centro, Veracruz, Veracruz, C.P. 91700.\\
\hline

\end{tabular}
\end{table}

\begin{table}[H]
\centering
\begin{tabular}{|c|m{.4\textwidth}|}
\hline
Obra& Dibujo	\\
\hline
Autor & Se atribuye a Ra\'ul Gonzalez Esquivel\\
\hline
T\'itulo & ``Conjunto escult\'orico. Sacerdotes Portadores de un Incensario. Cultura Totonaca. Centro de Veracruz (Mixtequilla) Cl\'asico Tard\'io, S. VI-IX D.C.''\\
\hline
Tem\'atica & Representaci\'on de arte precolombino, im\'agenes de deidades, figuras antropomorfas, cabezas colosales, culto religioso ind\'igena y personajes mitol\'ogicos.\\
\hline
T\'ecnicas &Grafito sobre papel \\
\hline
Medidas & 41.5 cm $\times$ 59 cm \\
\hline
 Formato & Peque\~no \\
 \hline
 Firma & Se atribuye a Ra\'ul Gonzalez Esquivel\\ 
 \hline
  Fecha & \\
 \hline
 Provenance & Colecci\'on particular\\
 \hline
 Certificado de autenticidad& Firmado por Pablo C. Goebel el 10 de abril de 1996.  \\
 \hline 
  Materiales empleados & Lapiz sobre papel\\
 \hline
 Caracter\'isticas especiales & Forma parte de una colecci\'on privada de 39 dibujos. 
Contiene dibujo con el texto ``Bufete Campos'' que corresponde a la obra por encargo. \\
\hline 
Ubicaci\'on & Calle Constituci\'on 124, colonia Centro, Veracruz, Veracruz, C.P. 91700.\\
\hline

\end{tabular}
\end{table}

\begin{table}[H]
\centering
\begin{tabular}{|c|m{.4\textwidth}|}
\hline
Obra& Dibujo	\\
\hline
Autor & Se atribuye a Ra\'ul Gonzalez Esquivel\\
\hline
T\'itulo & ``Figura Sonriente. Pintada y articulada. El Zapotal. Mpio. Ignacio de la Llave, Cl\'asico Tard\'io S. VI-IX D.C.''\\
\hline
Tem\'atica & Representaci\'on de arte precolombino, im\'agenes de deidades, figuras antropomorfas, cabezas colosales, culto religioso ind\'igena y personajes mitol\'ogicos.\\
\hline
T\'ecnicas &Grafito sobre papel \\
\hline
Medidas & 41.5 cm $\times$ 59 cm \\
\hline
 Formato & Peque\~no \\
 \hline
 Firma & Se atribuye a Ra\'ul Gonzalez Esquivel\\ 
 \hline
  Fecha & \\
 \hline
 Provenance & Colecci\'on particular\\
 \hline
 Certificado de autenticidad& Firmado por Pablo C. Goebel el 10 de abril de 1996.  \\
 \hline 
  Materiales empleados & Lapiz sobre papel\\
 \hline
 Caracter\'isticas especiales & Forma parte de una colecci\'on privada de 39 dibujos. 
Contiene dibujo con el texto ``Bufete Campos'' que corresponde a la obra por encargo. \\
\hline 
Ubicaci\'on & Calle Constituci\'on 124, colonia Centro, Veracruz, Veracruz, C.P. 91700.\\
\hline

\end{tabular}
\end{table}

\begin{table}[H]
\centering
\begin{tabular}{|c|m{.4\textwidth}|}
\hline
Obra& Dibujo	\\
\hline
Autor & Se atribuye a Ra\'ul Gonzalez Esquivel\\
\hline
T\'itulo &``Carita Sonriente. Cultura Totonaca. Los Cerros, Veracruz. Cl\'asico Tard\'io S. VI-IX D.C. Octubre 8 de 1993.'' \\
\hline
Tem\'atica & Representaci\'on de arte precolombino, im\'agenes de deidades, figuras antropomorfas, cabezas colosales, culto religioso ind\'igena y personajes mitol\'ogicos.\\
\hline
T\'ecnicas &Grafito sobre papel \\
\hline
Medidas & 41.5 cm $\times$ 59 cm \\
\hline
 Formato & Peque\~no \\
 \hline
 Firma & Se atribuye a Ra\'ul Gonzalez Esquivel\\ 
 \hline
  Fecha & Octubre 8 de 1993.\\
 \hline
 Provenance & Colecci\'on particular\\
 \hline
 Certificado de autenticidad& Firmado por Pablo C. Goebel el 10 de abril de 1996.  \\
 \hline 
  Materiales empleados & Lapiz sobre papel\\
 \hline
 Caracter\'isticas especiales & Forma parte de una colecci\'on privada de 39 dibujos. 
Contiene dibujo con el texto ``Bufete Campos'' que corresponde a la obra por encargo. \\
\hline 
Ubicaci\'on & Calle Constituci\'on 124, colonia Centro, Veracruz, Veracruz, C.P. 91700.\\
\hline

\end{tabular}
\end{table}

\begin{table}[H]
\centering
\begin{tabular}{|c|m{.4\textwidth}|}
\hline
Obra& Dibujo	\\
\hline
Autor & Se atribuye a Ra\'ul Gonzalez Esquivel\\
\hline
T\'itulo & ``Figura Danzante. Cultura Totonaca, Veracruz. Octubre 08 de 1993.''\\
\hline
Tem\'atica & Representaci\'on de arte precolombino, im\'agenes de deidades, figuras antropomorfas, cabezas colosales, culto religioso ind\'igena y personajes mitol\'ogicos.\\
\hline
T\'ecnicas &Grafito sobre papel \\
\hline
Medidas & 41.5 cm $\times$ 59 cm \\
\hline
 Formato & Peque\~no \\
 \hline
 Firma & Se atribuye a Ra\'ul Gonzalez Esquivel\\ 
 \hline
  Fecha & Octubre 8 de 1993.\\
 \hline
 Provenance & Colecci\'on particular\\
 \hline
 Certificado de autenticidad& Firmado por Pablo C. Goebel el 10 de abril de 1996.  \\
 \hline 
  Materiales empleados & Lapiz sobre papel\\
 \hline
 Caracter\'isticas especiales & Forma parte de una colecci\'on privada de 39 dibujos. 
Contiene dibujo con el texto ``Bufete Campos'' que corresponde a la obra por encargo. \\
\hline 
Ubicaci\'on & Calle Constituci\'on 124, colonia Centro, Veracruz, Veracruz, C.P. 91700.\\
\hline

\end{tabular}
\end{table}

\begin{table}[H]
\centering
\begin{tabular}{|c|m{.4\textwidth}|}
\hline
Obra& Dibujo	\\
\hline
Autor & Se atribuye a Ra\'ul Gonzalez Esquivel\\
\hline
T\'itulo &``Carita Sonriente. Cultura Totonaca S. VI-XI D.C. Los Cerros, Veracruz. Octubre 08 de 1993.'' \\
\hline
Tem\'atica & Representaci\'on de arte precolombino, im\'agenes de deidades, figuras antropomorfas, cabezas colosales, culto religioso ind\'igena y personajes mitol\'ogicos.\\
\hline
T\'ecnicas &Grafito sobre papel \\
\hline
Medidas & 41.5 cm $\times$ 59 cm \\
\hline
 Formato & Peque\~no \\
 \hline
 Firma & Se atribuye a Ra\'ul Gonzalez Esquivel\\ 
 \hline
  Fecha & Octubre 8 de 1993.\\
 \hline
 Provenance & Colecci\'on particular\\
 \hline
 Certificado de autenticidad& Firmado por Pablo C. Goebel el 10 de abril de 1996.  \\
 \hline 
  Materiales empleados & Lapiz sobre papel\\
 \hline
 Caracter\'isticas especiales & Forma parte de una colecci\'on privada de 39 dibujos. 
Contiene dibujo con el texto ``Bufete Campos'' que corresponde a la obra por encargo. \\
\hline 
Ubicaci\'on & Calle Constituci\'on 124, colonia Centro, Veracruz, Veracruz, C.P. 91700.\\
\hline

\end{tabular}
\end{table}

\begin{table}[H]
\centering
\begin{tabular}{|c|m{.4\textwidth}|}
\hline
Obra& Dibujo	\\
\hline
Autor & Se atribuye a Ra\'ul Gonzalez Esquivel\\
\hline
T\'itulo & ``Carita Sonriente. Cultura Totonaca. S. VI-XI D.C. Los Cerros, Veracruz, Octubre 08 de 1993.'' \\
\hline
Tem\'atica & Representaci\'on de arte precolombino, im\'agenes de deidades, figuras antropomorfas, cabezas colosales, culto religioso ind\'igena y personajes mitol\'ogicos.\\
\hline
T\'ecnicas &Grafito sobre papel \\
\hline
Medidas & 41.5 cm $\times$ 59 cm \\
\hline
 Formato & Peque\~no \\
 \hline
 Firma & Se atribuye a Ra\'ul Gonzalez Esquivel\\ 
 \hline
  Fecha & Octubre 8 de 1993.\\
 \hline
 Provenance & Colecci\'on particular\\
 \hline
 Certificado de autenticidad& Firmado por Pablo C. Goebel el 10 de abril de 1996.  \\
 \hline 
  Materiales empleados & Lapiz sobre papel\\
 \hline
 Caracter\'isticas especiales & Forma parte de una colecci\'on privada de 39 dibujos. 
Contiene dibujo con el texto ``Bufete Campos'' que corresponde a la obra por encargo. \\
\hline 
Ubicaci\'on & Calle Constituci\'on 124, colonia Centro, Veracruz, Veracruz, C.P. 91700.\\
\hline

\end{tabular}
\end{table}

\begin{table}[H]
\centering
\begin{tabular}{|c|m{.4\textwidth}|}
\hline
Obra& Dibujo	\\
\hline
Autor & Se atribuye a Ra\'ul Gonzalez Esquivel\\
\hline
T\'itulo & ``Figura SOnriente Desnuda. El Zapotal. Mpio. de Ignacio de la Llave. Veracruz. Cultura Totonaca. Cl\'asico tard\'io. S. VI-IX D.C. Octubre 3 de 1993.''\\
\hline
Tem\'atica & Representaci\'on de arte precolombino, im\'agenes de deidades, figuras antropomorfas, cabezas colosales, culto religioso ind\'igena y personajes mitol\'ogicos.\\
\hline
T\'ecnicas &Grafito sobre papel \\
\hline
Medidas & 41.5 cm $\times$ 59 cm \\
\hline
 Formato & Peque\~no \\
 \hline
 Firma & Se atribuye a Ra\'ul Gonzalez Esquivel\\ 
 \hline
  Fecha & Octubree 3 de 1993.\\
 \hline
 Provenance & Colecci\'on particular\\
 \hline
 Certificado de autenticidad& Firmado por Pablo C. Goebel el 10 de abril de 1996.  \\
 \hline 
  Materiales empleados & Lapiz sobre papel\\
 \hline
 Caracter\'isticas especiales & Forma parte de una colecci\'on privada de 39 dibujos. 
Contiene dibujo con el texto ``Bufete Campos'' que corresponde a la obra por encargo. \\
\hline 
Ubicaci\'on & Calle Constituci\'on 124, colonia Centro, Veracruz, Veracruz, C.P. 91700.\\
\hline

\end{tabular}
\end{table}

\begin{table}[H]
\centering
\begin{tabular}{|c|m{.4\textwidth}|}
\hline
Obra& Dibujo	\\
\hline
Autor & Se atribuye a Ra\'ul Gonzalez Esquivel\\
\hline
T\'itulo & ``Figura S\'edente. Cultura Totonaca. El Salto. Mpio de Ignacio de la Llave. Veracruz. Octubre 4 de 1993.'' \\
\hline
Tem\'atica & Representaci\'on de arte precolombino, im\'agenes de deidades, figuras antropomorfas, cabezas colosales, culto religioso ind\'igena y personajes mitol\'ogicos.\\
\hline
T\'ecnicas &Grafito sobre papel \\
\hline
Medidas & 41.5 cm $\times$ 59 cm \\
\hline
 Formato & Peque\~no \\
 \hline
 Firma & Se atribuye a Ra\'ul Gonzalez Esquivel\\ 
 \hline
  Fecha & Octubre 4 de 1993. \\
 \hline
 Provenance & Colecci\'on particular\\
 \hline
 Certificado de autenticidad& Firmado por Pablo C. Goebel el 10 de abril de 1996.  \\
 \hline 
  Materiales empleados & Lapiz sobre papel\\
 \hline
 Caracter\'isticas especiales & Forma parte de una colecci\'on privada de 39 dibujos. 
Contiene dibujo con el texto ``Bufete Campos'' que corresponde a la obra por encargo. \\
\hline 
Ubicaci\'on & Calle Constituci\'on 124, colonia Centro, Veracruz, Veracruz, C.P. 91700.\\
\hline

\end{tabular}
\end{table}

\begin{table}[H]
\centering
\begin{tabular}{|c|m{.4\textwidth}|}
\hline
Obra& Dibujo	\\
\hline
Autor & Se atribuye a Ra\'ul Gonzalez Esquivel\\
\hline
T\'itulo & ``Figura Sonriente. Cultura Totonaca. Veracruz. Octubre 2 de 1993.'' \\
\hline
Tem\'atica & Representaci\'on de arte precolombino, im\'agenes de deidades, figuras antropomorfas, cabezas colosales, culto religioso ind\'igena y personajes mitol\'ogicos.\\
\hline
T\'ecnicas &Grafito sobre papel \\
\hline
Medidas & 41.5 cm $\times$ 59 cm \\
\hline
 Formato & Peque\~no \\
 \hline
 Firma & Se atribuye a Ra\'ul Gonzalez Esquivel\\ 
 \hline
  Fecha & Octubre 2 de 1993.\\
 \hline
 Provenance & Colecci\'on particular\\
 \hline
 Certificado de autenticidad& Firmado por Pablo C. Goebel el 10 de abril de 1996.  \\
 \hline 
  Materiales empleados & Lapiz sobre papel\\
 \hline
 Caracter\'isticas especiales & Forma parte de una colecci\'on privada de 39 dibujos. 
Contiene dibujo con el texto ``Bufete Campos'' que corresponde a la obra por encargo. \\
\hline 
Ubicaci\'on & Calle Constituci\'on 124, colonia Centro, Veracruz, Veracruz, C.P. 91700.\\
\hline

\end{tabular}
\end{table}

\begin{table}[H]
\centering
\begin{tabular}{|c|m{.4\textwidth}|}
\hline
Obra& Dibujo	\\
\hline
Autor & Se atribuye a Ra\'ul Gonzalez Esquivel\\
\hline
T\'itulo & ''Hacha Votiva con cara de guerrero. Diciembre 10 de 1993.'' \\
\hline
Tem\'atica & Representaci\'on de arte precolombino, im\'agenes de deidades, figuras antropomorfas, cabezas colosales, culto religioso ind\'igena y personajes mitol\'ogicos.\\
\hline
T\'ecnicas &Grafito sobre papel \\
\hline
Medidas & 41.5 cm $\times$ 59 cm \\
\hline
 Formato & Peque\~no \\
 \hline
 Firma & Se atribuye a Ra\'ul Gonzalez Esquivel\\ 
 \hline
  Fecha & Diciembre 10 de 1993. \\
 \hline
 Provenance & Colecci\'on particular\\
 \hline
 Certificado de autenticidad& Firmado por Pablo C. Goebel el 10 de abril de 1996.  \\
 \hline 
  Materiales empleados & Lapiz sobre papel\\
 \hline
 Caracter\'isticas especiales & Forma parte de una colecci\'on privada de 39 dibujos. 
Contiene dibujo con el texto ``Bufete Campos'' que corresponde a la obra por encargo. \\
\hline 
Ubicaci\'on & Calle Constituci\'on 124, colonia Centro, Veracruz, Veracruz, C.P. 91700.\\
\hline

\end{tabular}
\end{table}

\begin{table}[H]
\centering
\begin{tabular}{|c|m{.4\textwidth}|}
\hline
Obra& Dibujo	\\
\hline
Autor & Se atribuye a Ra\'ul Gonzalez Esquivel\\
\hline
T\'itulo & ``Hacha Votiva, representa a un guerrero muerto. Napatecuhtlan, Mpio. de Perote. Cl\'asico Tard\'io S. VI-IX. Diciembre 10 de 1993.'' \\
\hline
Tem\'atica & Representaci\'on de arte precolombino, im\'agenes de deidades, figuras antropomorfas, cabezas colosales, culto religioso ind\'igena y personajes mitol\'ogicos.\\
\hline
T\'ecnicas &Grafito sobre papel \\
\hline
Medidas & 41.5 cm $\times$ 59 cm \\
\hline
 Formato & Peque\~no \\
 \hline
 Firma & Se atribuye a Ra\'ul Gonzalez Esquivel\\ 
 \hline
  Fecha & Diciembre 10 de 1993.\\
 \hline
 Provenance & Colecci\'on particular\\
 \hline
 Certificado de autenticidad& Firmado por Pablo C. Goebel el 10 de abril de 1996.  \\
 \hline 
  Materiales empleados & Lapiz sobre papel\\
 \hline
 Caracter\'isticas especiales & Forma parte de una colecci\'on privada de 39 dibujos. 
Contiene dibujo con el texto ``Bufete Campos'' que corresponde a la obra por encargo. \\
\hline 
Ubicaci\'on & Calle Constituci\'on 124, colonia Centro, Veracruz, Veracruz, C.P. 91700.\\
\hline

\end{tabular}
\end{table}

\begin{table}[H]
\centering
\begin{tabular}{|c|m{.4\textwidth}|}
\hline
Obra& Dibujo	\\
\hline
Autor & Se atribuye a Ra\'ul Gonzalez Esquivel\\
\hline
T\'itulo & ``Tlasolteotl. Santana, Mpio. de Tlalixcoyan. Cl\'asico Tard\'io S. VI-IX D.C. Noviembre 26 de 1993.''\\
\hline
Tem\'atica & Representaci\'on de arte precolombino, im\'agenes de deidades, figuras antropomorfas, cabezas colosales, culto religioso ind\'igena y personajes mitol\'ogicos.\\
\hline
T\'ecnicas &Grafito sobre papel \\
\hline
Medidas & 41.5 cm $\times$ 59 cm \\
\hline
 Formato & Peque\~no \\
 \hline
 Firma & Se atribuye a Ra\'ul Gonzalez Esquivel\\ 
 \hline
  Fecha & Noviembre 26 de 1993. \\
 \hline
 Provenance & Colecci\'on particular\\
 \hline
 Certificado de autenticidad& Firmado por Pablo C. Goebel el 10 de abril de 1996.  \\
 \hline 
  Materiales empleados & Lapiz sobre papel\\
 \hline
 Caracter\'isticas especiales & Forma parte de una colecci\'on privada de 39 dibujos. 
Contiene dibujo con el texto ``Bufete Campos'' que corresponde a la obra por encargo. \\
\hline 
Ubicaci\'on & Calle Constituci\'on 124, colonia Centro, Veracruz, Veracruz, C.P. 91700.\\
\hline

\end{tabular}
\end{table}

\begin{table}[H]
\centering
\begin{tabular}{|c|m{.4\textwidth}|}
\hline
Obra& Dibujo	\\
\hline
Autor & Se atribuye a Ra\'ul Gonzalez Esquivel\\
\hline
T\'itulo & ``Cinuateteotl no. 2. El Cocuhite, Mpio. de Tlalixcoyan. Cl\'asico Tard\'io S. VI-IX D.C. Noviembre 27 de 1993.''\\
\hline
Tem\'atica & Representaci\'on de arte precolombino, im\'agenes de deidades, figuras antropomorfas, cabezas colosales, culto religioso ind\'igena y personajes mitol\'ogicos.\\
\hline
T\'ecnicas &Grafito sobre papel \\
\hline
Medidas & 41.5 cm $\times$ 59 cm \\
\hline
 Formato & Peque\~no \\
 \hline
 Firma & Se atribuye a Ra\'ul Gonzalez Esquivel\\ 
 \hline
  Fecha & Noviembre 27 de 1993.\\
 \hline
 Provenance & Colecci\'on particular\\
 \hline
 Certificado de autenticidad& Firmado por Pablo C. Goebel el 10 de abril de 1996.  \\
 \hline 
  Materiales empleados & Lapiz sobre papel\\
 \hline
 Caracter\'isticas especiales & Forma parte de una colecci\'on privada de 39 dibujos. 
Contiene dibujo con el texto ``Bufete Campos'' que corresponde a la obra por encargo. \\
\hline 
Ubicaci\'on & Calle Constituci\'on 124, colonia Centro, Veracruz, Veracruz, C.P. 91700.\\
\hline

\end{tabular}
\end{table}

\begin{table}[H]
\centering
\begin{tabular}{|c|m{.4\textwidth}|}
\hline
Obra& Dibujo	\\
\hline
Autor & Se atribuye a Ra\'ul Gonzalez Esquivel\\
\hline
T\'itulo & Cihuateteotl no.1, del Cocuhite, Mpio. de Tlalixcoyan. Cl\'asico Tard\'io. S. VI-IX. Noviembre 18 de 1998.''\\
\hline
Tem\'atica & Representaci\'on de arte precolombino, im\'agenes de deidades, figuras antropomorfas, cabezas colosales, culto religioso ind\'igena y personajes mitol\'ogicos.\\
\hline
T\'ecnicas &Grafito sobre papel \\
\hline
Medidas & 41.5 cm $\times$ 59 cm \\
\hline
 Formato & Peque\~no \\
 \hline
 Firma & Se atribuye a Ra\'ul Gonzalez Esquivel\\ 
 \hline
  Fecha & Noviembre 18 de 1993.\\
 \hline
 Provenance & Colecci\'on particular\\
 \hline
 Certificado de autenticidad& Firmado por Pablo C. Goebel el 10 de abril de 1996.  \\
 \hline 
  Materiales empleados & Lapiz sobre papel\\
 \hline
 Caracter\'isticas especiales & Forma parte de una colecci\'on privada de 39 dibujos. 
Contiene dibujo con el texto ``Bufete Campos'' que corresponde a la obra por encargo. \\
\hline 
Ubicaci\'on & Calle Constituci\'on 124, colonia Centro, Veracruz, Veracruz, C.P. 91700.\\
\hline

\end{tabular}
\end{table}

\begin{table}[H]
\centering
\begin{tabular}{|c|m{.4\textwidth}|}
\hline
Obra& Dibujo	\\
\hline
Autor & Se atribuye a Ra\'ul Gonzalez Esquivel\\
\hline
T\'itulo & ``Danzante. Centro de Veracruz. Cl\'asico Tard\'io S. VI-IX. Noviembre 10 de 1993.'' \\
\hline
Tem\'atica & Representaci\'on de arte precolombino, im\'agenes de deidades, figuras antropomorfas, cabezas colosales, culto religioso ind\'igena y personajes mitol\'ogicos.\\
\hline
T\'ecnicas &Grafito sobre papel \\
\hline
Medidas & 41.5 cm $\times$ 59 cm \\
\hline
 Formato & Peque\~no \\
 \hline
 Firma & Se atribuye a Ra\'ul Gonzalez Esquivel\\ 
 \hline
  Fecha & Noviembre 10 de 1993.\\
 \hline
 Provenance & Colecci\'on particular\\
 \hline
 Certificado de autenticidad& Firmado por Pablo C. Goebel el 10 de abril de 1996.  \\
 \hline 
  Materiales empleados & Lapiz sobre papel\\
 \hline
 Caracter\'isticas especiales & Forma parte de una colecci\'on privada de 39 dibujos. 
Contiene dibujo con el texto ``Bufete Campos'' que corresponde a la obra por encargo. \\
\hline 
Ubicaci\'on & Calle Constituci\'on 124, colonia Centro, Veracruz, Veracruz, C.P. 91700.\\
\hline

\end{tabular}
\end{table}

\begin{table}[H]
\centering
\begin{tabular}{|c|m{.4\textwidth}|}
\hline
Obra& Dibujo	\\
\hline
Autor & Se atribuye a Ra\'ul Gonzalez Esquivel\\
\hline
T\'itulo & ``Figura Sonriente Femenina. El Zapotal, Mpio. de Ignacio de la Llave. Cl\'asico Tard\'io S. VI-IX. Noviembre 13 de 1993.''\\
\hline
Tem\'atica & Representaci\'on de arte precolombino, im\'agenes de deidades, figuras antropomorfas, cabezas colosales, culto religioso ind\'igena y personajes mitol\'ogicos.\\
\hline
T\'ecnicas &Grafito sobre papel \\
\hline
Medidas & 41.5 cm $\times$ 59 cm \\
\hline
 Formato & Peque\~no \\
 \hline
 Firma & Se atribuye a Ra\'ul Gonzalez Esquivel\\ 
 \hline
  Fecha & Noviembre 13 de 1993. \\
 \hline
 Provenance & Colecci\'on particular\\
 \hline
 Certificado de autenticidad& Firmado por Pablo C. Goebel el 10 de abril de 1996.  \\
 \hline 
  Materiales empleados & Lapiz sobre papel\\
 \hline
 Caracter\'isticas especiales & Forma parte de una colecci\'on privada de 39 dibujos. 
Contiene dibujo con el texto ``Bufete Campos'' que corresponde a la obra por encargo. \\
\hline 
Ubicaci\'on & Calle Constituci\'on 124, colonia Centro, Veracruz, Veracruz, C.P. 91700.\\
\hline

\end{tabular}
\end{table}

\begin{table}[H]
\centering
\begin{tabular}{|c|m{.4\textwidth}|}
\hline
Obra& Dibujo	\\
\hline
Autor & Se atribuye a Ra\'ul Gonzalez Esquivel\\
\hline
T\'itulo & ``Jugador de Pelota. Mata Redonda, Mpio. de Villa Cuauht\'emoc. S. II-IX D.C. Periodo Panuco. Noviembre 1 de 1993.''\\
\hline
Tem\'atica & Representaci\'on de arte precolombino, im\'agenes de deidades, figuras antropomorfas, cabezas colosales, culto religioso ind\'igena y personajes mitol\'ogicos.\\
\hline
T\'ecnicas &Grafito sobre papel \\
\hline
Medidas & 41.5 cm $\times$ 59 cm \\
\hline
 Formato & Peque\~no \\
 \hline
 Firma & Se atribuye a Ra\'ul Gonzalez Esquivel\\ 
 \hline
  Fecha & Noviembre 1 de 1993.\\
 \hline
 Provenance & Colecci\'on particular\\
 \hline
 Certificado de autenticidad& Firmado por Pablo C. Goebel el 10 de abril de 1996.  \\
 \hline 
  Materiales empleados & Lapiz sobre papel\\
 \hline
 Caracter\'isticas especiales & Forma parte de una colecci\'on privada de 39 dibujos. 
Contiene dibujo con el texto ``Bufete Campos'' que corresponde a la obra por encargo. \\
\hline 
Ubicaci\'on & Calle Constituci\'on 124, colonia Centro, Veracruz, Veracruz, C.P. 91700.\\
\hline

\end{tabular}
\end{table}

\begin{table}[H]
\centering
\begin{tabular}{|c|m{.4\textwidth}|}
\hline
Obra& Dibujo	\\
\hline
Autor & Se atribuye a Ra\'ul Gonzalez Esquivel\\
\hline
T\'itulo & ``Tlasolteotl, con Gran Quechquemitl  Napiola, Mpio. de Tierra Blanca, Ver. Noviembre 03 de 1993.''\\
\hline
Tem\'atica & Representaci\'on de arte precolombino, im\'agenes de deidades, figuras antropomorfas, cabezas colosales, culto religioso ind\'igena y personajes mitol\'ogicos.\\
\hline
T\'ecnicas &Grafito sobre papel \\
\hline
Medidas & 41.5 cm $\times$ 59 cm \\
\hline
 Formato & Peque\~no \\
 \hline
 Firma & Se atribuye a Ra\'ul Gonzalez Esquivel\\ 
 \hline
  Fecha & Noviembre 3 de 1993.\\
 \hline
 Provenance & Colecci\'on particular\\
 \hline
 Certificado de autenticidad& Firmado por Pablo C. Goebel el 10 de abril de 1996.  \\
 \hline 
  Materiales empleados & Lapiz sobre papel\\
 \hline
 Caracter\'isticas especiales & Forma parte de una colecci\'on privada de 39 dibujos. 
Contiene dibujo con el texto ``Bufete Campos'' que corresponde a la obra por encargo. \\
\hline 
Ubicaci\'on & Calle Constituci\'on 124, colonia Centro, Veracruz, Veracruz, C.P. 91700.\\
\hline

\end{tabular}
\end{table}

\begin{table}[H]
\centering
\begin{tabular}{|c|m{.4\textwidth}|}
\hline
Obra& Dibujo	\\
\hline
Autor & Se atribuye a Ra\'ul Gonzalez Esquivel\\
\hline
T\'itulo & ``Hacha Votiva en forma de garza. Mpio de Coxquihui, VEracruz. Cla\'asico Tard\'io S. VI-IX D.C. Diciembre 10 de 1993.''\\
\hline
Tem\'atica & Representaci\'on de arte precolombino, im\'agenes de deidades, figuras antropomorfas, cabezas colosales, culto religioso ind\'igena y personajes mitol\'ogicos.\\
\hline
T\'ecnicas &Grafito sobre papel \\
\hline
Medidas & 41.5 cm $\times$ 59 cm \\
\hline
 Formato & Peque\~no \\
 \hline
 Firma & Se atribuye a Ra\'ul Gonzalez Esquivel\\ 
 \hline
  Fecha & Diciembre 10 de 1993.\\
 \hline
 Provenance & Colecci\'on particular\\
 \hline
 Certificado de autenticidad& Firmado por Pablo C. Goebel el 10 de abril de 1996.  \\
 \hline 
  Materiales empleados & Lapiz sobre papel\\
 \hline
 Caracter\'isticas especiales & Forma parte de una colecci\'on privada de 39 dibujos. 
Contiene dibujo con el texto ``Bufete Campos'' que corresponde a la obra por encargo. \\
\hline 
Ubicaci\'on & Calle Constituci\'on 124, colonia Centro, Veracruz, Veracruz, C.P. 91700.\\
\hline

\end{tabular}
\end{table}


\end{document}
